\documentclass[11pt,a4paper]{article}
\usepackage{ConfigMP}

% ================================
% INFORMACIÓN DEL DOCUMENTO
% ================================
% Sustituya los valores por la información correspondiente

% Información del miniproyecto
\newcommand{\claveproyecto}{MP-26-XXX}
\newcommand{\tituloproyecto}{Escriba un titulo claro, especifico y representativo del trabajo realizado}
\newcommand{\autores}{Incluya los nombres completos en el orden de participacion y use superindices}
\newcommand{\adscripciones}{Indique departamento o programa, institucion}

\begin{document}

% ================================
% TÍTULO Y INFORMACIÓN DEL PROYECTO
% ================================
\doctitle{Sustituya los textos entre corchetes y elimine las orientaciones antes de entregar el reporte final.}

\projectdata{\claveproyecto}{\tituloproyecto}{\autores}{\adscripciones}

% ================================
% INTRODUCCIÓN
% ================================
\section{Introduccion}

\begin{sectioncontent}
Presente el contexto, los antecedentes esenciales, la relevancia y la justificacion del miniproyecto. Cierre con el proposito y el alcance del trabajo e incorpore las citas correspondientes.
\end{sectioncontent}

\vspace{10pt}

% ================================
% OBJETIVOS
% ================================
\section{Objetivos}

\begin{sectioncontent}
Formule un objetivo general y, cuando corresponda, objetivos especificos con verbos en infinitivo. Deben ser claros, verificables y congruentes con la metodologia y los resultados.

\subsection*{Objetivo General}
Escriba aqui el objetivo general

\subsection*{Objetivos Especificos}
\begin{itemize}
    \item Objetivo especifico 1
    \item Objetivo especifico 2
    \item Objetivo especifico 3
\end{itemize}
\end{sectioncontent}

\vspace{10pt}

% ================================
% MATERIALES Y MÉTODOS
% ================================
\section{Materiales y metodos}

\begin{sectioncontent}
Describa materiales, reactivos, equipos, diseño experimental y procedimientos con detalle reproducible. Incluya variables, condiciones, analisis de datos, normas y aspectos de seguridad o etica aplicables; use subsecciones si resulta util.

\subsection*{Materiales}
\begin{itemize}
    \item Material 1
    \item Material 2
    \item Material 3
\end{itemize}

\subsection*{Equipos}
\begin{itemize}
    \item Equipo 1
    \item Equipo 2
    \item Equipo 3
\end{itemize}

\subsection*{Procedimiento}
\begin{enumerate}
    \item Paso 1
    \item Paso 2
    \item Paso 3
\end{enumerate}
\end{sectioncontent}

\vspace{10pt}

% ================================
% RESULTADOS Y DISCUSIONES
% ================================
\section{Resultados y discusiones}

\begin{sectioncontent}
Presente los hallazgos de forma objetiva y en el orden de los objetivos. Use tablas y figuras numeradas, titulos o pies explicativos, unidades y analisis estadistico cuando corresponda, y cite cada elemento en el texto.

% Ejemplo de tabla
\begin{table}[H]
    \centering
    \caption{Titulo descriptivo de la tabla}
    \label{tab:ejemplo}
    \begin{tabular}{@{}lcc@{}}
        \toprule
        \textbf{Variable} & \textbf{Valor 1} & \textbf{Valor 2} \\
        \midrule
        Resultado 1 & 0.00 & 0.00 \\
        Resultado 2 & 0.00 & 0.00 \\
        Resultado 3 & 0.00 & 0.00 \\
        \bottomrule
    \end{tabular}
\end{table}

\subsection*{Discusion}
Analice e interprete los resultados obtenidos, compare con la literatura existente y discuta las implicaciones de los hallazgos.
\end{sectioncontent}

\vspace{10pt}

% ================================
% CONCLUSIONES
% ================================
\section{Conclusiones}

\begin{sectioncontent}
Sintetice los resultados principales y señale en que medida se cumplieron los objetivos. No introduzca informacion nueva; mencione limitaciones, implicaciones o trabajo futuro solo cuando sean pertinentes.

\begin{itemize}
    \item Conclusion 1
    \item Conclusion 2
    \item Conclusion 3
\end{itemize}
\end{sectioncontent}

\vspace{10pt}

% ================================
% REFERENCIAS
% ================================
\section{Referencias}

\begin{orientation}
    \textbf{Orientacion:} Incluya unicamente las fuentes citadas y verifique su correspondencia con la lista final. Mantenga un solo estilo bibliografico e incorpore DOI o URL cuando proceda.
\end{orientation}

\vspace{8pt}

% Ejemplo de formato de referencias (estilo APA)
[1] Apellido, A. A. (2026). \textit{Titulo del articulo}. Nombre de la revista, \textit{Volumen}(Numero), paginas. DOI o URL

[2] Apellido, B. B. (2026). \textit{Titulo del libro}. Editorial.

[3] Apellido, D. D. (2026). Titulo del capitulo. En E. E. Apellido (Ed.), \textit{Titulo del libro} (pp. xx-xx). Editorial.

\end{document}